% Both A authors. Local-only tests must not be presented as full benchmarks.
\subsection{Supporting experiments, validation and final selection}
An unrun proposal, a13, would have imposed a fixed lock order on T7. We did
not run it because the preserved baseline recorded serialization failures,
not deadlock failures. A deadlock-prevention rewrite would not directly
address that error category \citep{pg-isolation}.

Local-only implementations explored a one-call T1 function, maintained
summaries, memory settings, executor concurrency and alternative read-only
declarations. They supported implementation and correctness testing; their
small datasets and short runs do not substitute for course-cluster benchmarks.
Their code is packaged under \code{supporting_implementations/local_only/},
with outputs under \code{evidence/local-only/}.

Validation covered inventory non-negativity, order totals, identifier uniqueness,
rollback behaviour, placement metadata and selected transaction semantics.
Tests included transfers with different inventory and product identifiers.
Validation ran outside the timed interval. Passing checks establish the tested
invariants at inspected checkpoints, not correctness under every concurrent
schedule. Passing checks after an abort do not make a performance run successful.

\textbf{A112 is the final Workload A design.} It supplies a complete successful
seed-42 matrix and a substantial matched improvement. A115 is a promising but
unreliable follow-up. A10 demonstrates a separate transaction-level improvement
that has not yet been combined with a112.
